\chapter{Inflection}

\begin{sourcebox}
A Russian noun will have a root, indicating its basic meaning, and a case-ending, indicating its relationship to the rest of the sentence; and it is said to be \emph{declined}, or \emph{inflected}, when all its possible case-forms are displayed in a table. Prepositions can be used not only with various cases of nouns but also as prefixes to verbs, e.g. \emph{circumscribe, describe, inscribe}, etc. This second role of prepositions is fundamental for the study of Russian vocabulary.
\end{sourcebox}

\section{The concept of declension}

The Greek philosopher Aristotle (384--322 B.C.) was the first, in the Western world, to make a systematic study of grammar, and the numerous grammatical terms still in use in our schools are either the original Greek words (\emph{phrase, comma, colon}, etc.) standardized by him and his Alexandrian successors, or else are Latin translations of such words (\emph{adjective, verb, neuter, case}, etc.). His statements about Greek grammar are illuminating for the grammar of other Indo-European languages as well, including English and Russian.

Noting that a Greek noun (the name of a person or thing) changes its ending according to whether it represents: i) the subject performing the action, ii) an entity closely connected with the subject, iii) the indirect object of the action, or iv) the direct object, the ancient grammarians, taking their terminology from the various possible positions for the fall of (non-cubical) Greek dice, described the first of the four situations as the direct or upright case (``case'' is the Latin word for fall, cf. \emph{casualty}, etc.) while the other cases were successively less upright; and the four cases, in that order, were said to decline from upright to horizontal. Those parts of speech that show case-endings, i.e. the nouns, pronouns and adjectives, are said to undergo \emph{declension} (i.e. falling), but a verb is said to be \emph{conjugated} (i.e. its forms are conjoined into a table) and the word \emph{inflection} (change of ending) is used for both declension and conjugation. Adverbs, prepositions and conjunctions, the other three of the traditional seven parts of speech, are uninflected.

The first of the four cases was called the \emph{nominative} (used to name the subject); the second (usually translated by \emph{of} in English) was the \emph{genitive} (from the Latin \emph{genus}, as indicating the class or entity to which the subject belonged; e.g. a member of the senate, the hand of Caesar, etc.); the third was the \emph{dative} (usually translated by \emph{to} or \emph{for} in English), from the Latin \emph{dare} (to give); e.g., ``he gave the boy a book'', or ``he gave a book to the boy''; ``he made the boy a kite'', or ``he made a kite for the boy''; and the Greek name for the fourth case, indicating the object of the action, was carelessly mistranslated into Latin as the \emph{accusative} (a better name would be \emph{objective}). In present-day Russian there are two more cases (also inherited from Indo-European), namely the \emph{instrumental} and the \emph{locative}, showing the means by which and the place where the action is performed. The locative case is used only after certain prepositions, for which reason it is also called the prepositional.

Verbs which have a direct object (i.e. in the accusative case) are called transitive and may or may not have an indirect object as well (dative case). In English a verb cannot have an indirect object unless it also has a direct object (i.e. is transitive) but in Russian (and in many other languages) some intransitive verbs may have an indirect object; in a sentence like \emph{the kite satisfies the boy} (or \emph{the function satisfies the condition}) the Russian word for \emph{boy} (or \emph{condition}) will be in the dative case, the verb \emph{satisfies} having some such force as \emph{is sufficient for}. Such verbs are said to ``be followed by the dative'' or to ``take the dative''. The dative is also very common after adjectives like \emph{equal to, opposite to}, etc.

In the same way certain verbs ``take the instrumental''. In the sentence \emph{let us use these results} the Russian word for \emph{results} will be in the instrumental case, since the Russian verb for \emph{use} has some such force as ``help ourselves by means of''. A few verbs ``take the genitive''; for example, in a sentence like \emph{the line touches the circle} the Russian word for \emph{circle} will be in the genitive case, suggesting that only part of the circle is touched. In a negative sentence (suggesting no part of) the Russian accusative after a transitive verb may be replaced by the genitive; in \emph{he did not make a kite} the Russian word for \emph{kite} may be either in the accusative or in the genitive.

\section{The three declensions}

The ancient grammarians also used the word \emph{declension} in a more concrete sense, namely to mean a definite set of endings indicating the various cases of a given noun. In modern Russian there are three such declensions, i.e. sets of case-endings, the first two of which it is convenient to examine together. Note that here we give the endings only for the singular, since in the plural (see §4) the three declensions coincide almost completely.

\begin{center}
\small
\begin{tabular}{@{}lcccccc@{}}
\toprule
& nom. & gen. & dat. & acc. & instr. & loc.\\
\midrule
1st declension & --/\ru{о} & \ru{а} & \ru{у} & --/\ru{о} & \ru{ом} & \ru{е}\\
2nd declension & \ru{а} & \ru{ы} & \ru{е} & \ru{у} & \ru{ой} & \ru{е}\\
\bottomrule
\end{tabular}
\end{center}

Here the symbol --/\ru{о} means that some nouns of the first declension, e.g. \ru{факт} \emph{fact}, have no overt ending in the nominative case, while others, e.g. \ru{место} \emph{place, locus}, have the ending \ru{-о}. For the reasons given below in §6, words like \ru{факт} are masculine, words like \ru{место} are neuter, and words of the second declension, e.g. \ru{сумма} \emph{sum}, are feminine.

Thus the three words \ru{факт}, \ru{место} and \ru{сумма} are declined in the singular as follows:

\begingroup
\small
\setlength{\tabcolsep}{2pt}
\renewcommand{\arraystretch}{1.06}
\begin{longtable}{@{}p{9mm} >{\cyrillicfont}p{19mm} p{23mm} >{\cyrillicfont}p{19mm} p{23mm} >{\cyrillicfont}p{19mm} p{23mm}@{}}
\toprule
& \multicolumn{2}{c}{1st declension: masculine} & \multicolumn{2}{c}{1st declension: neuter} & \multicolumn{2}{c}{2nd declension: feminine}\\
\midrule
nom. & факт & fact & место & place, locus & сумма & sum\\
gen. & факта & of a fact & места & of a locus & суммы & of a sum\\
dat. & факту & to a fact & месту & to a locus & сумме & to a sum\\
acc. & факт & fact & место & locus & сумму & sum\\
instr. & фактом & by means of a fact & местом & by means of a locus & суммой & by means of a sum\\
loc. & (о) факте & (about) a fact & (о) месте & (about) a locus & (о) сумме & (about) a sum\\
\bottomrule
\end{longtable}
\endgroup

Note that if the stem of a feminine noun ends in \ru{к, г, х} or a sibilant, the spelling rule (Chap. I, §6) demands that the ending \ru{ы} be replaced by \ru{и}, e.g. \ru{точка, точки} ``point.''

Note also that there are no definite or indefinite articles in Russian: only context decides whether \ru{факт} means \emph{fact}, \emph{the fact}, or \emph{a fact}. In reciting the declension of a noun aloud it is customary to include with the prepositional/locative the preposition \ru{о} (\ru{об} before many vowel-initial forms and \ru{обо} in a limited set of combinations).

For greater ease of pronunciation, some nouns of the \ru{факт} declension, e.g. \ru{коне́ц} \emph{end}, \ru{отре́зок} \emph{segment}, insert a so-called movable vowel \ru{е} (or \ru{о} before the back consonants \ru{к, г, х}) in the nominative and accusative singular, where the word would otherwise end in two consonants. Thus \ru{коне́ц, конца́} (stem \ru{конц-}) \emph{end}, \ru{отре́зок, отре́зка} (stem \ru{отрезк-}) \emph{segment}.

In the nouns \ru{факт} and \ru{место} the stem ends in a hard consonant. For corresponding nouns with soft stems, e.g. \ru{числи́тель} \emph{numerator} (masc.) and \ru{по́ле} \emph{field} (neut.), the case spellings use the front-vowel series; schematically:
\begin{center}
\ru{--/е, я, ю, --/е, ем, е}.
\end{center}

\begin{center}
\small
\begin{tabular}{@{}l >{\cyrillicfont}l >{\cyrillicfont}l@{}}
\toprule
& masculine & neuter\\
\midrule
nom. & числитель & поле\\
gen. & числителя & поля\\
dat. & числителю & полю\\
acc. & числитель & поле\\
instr. & числителем & полем\\
loc. & (о) числителе & (о) поле\\
\bottomrule
\end{tabular}
\end{center}

Finally, neuter and feminine stems may also end in the sequence written \ru{-и-} before their case endings. Thus \ru{сложе́ние} \emph{addition} (neut.) and \ru{фу́нкция} \emph{function} (fem.) are declined as follows (note \ru{-ии} in the prepositional of nouns in \ru{-ия} and in the dative feminine):

\begin{center}
\small
\begin{tabular}{@{}l >{\cyrillicfont}l >{\cyrillicfont}l@{}}
\toprule
& neuter & feminine\\
\midrule
nom. & сложение & функция\\
gen. & сложения & функции\\
dat. & сложению & функции\\
acc. & сложение & функцию\\
instr. & сложением & функцией\\
loc. & (о) сложении & (о) функции\\
\bottomrule
\end{tabular}
\end{center}

There remains the third declension, whose regular members are feminine. Note that several of the forms have the ending \ru{-и}:

\begin{center}
\small
\begin{tabular}{@{}l >{\cyrillicfont}l l@{}}
nom. & коммутативность & commutativity\\
gen. & коммутативности & \\
dat. & коммутативности & \\
acc. & коммутативность & \\
instr. & коммутативностью & \\
loc. & (о) коммутативности & \\
\end{tabular}
\end{center}


\section{Frequency of occurrence of nouns of the eight types}

The two classes of nouns commonest in mathematics are the ones represented by \ru{факт} and \ru{сложе́ние}. The \ru{факт} class includes genuine old Russian words like \ru{вид} \emph{form}, \ru{знак} \emph{sign}, and also many transliterations of foreign words like \ru{дифференциа́л, идеа́л, класс, коэффицие́нт, фо́кус, элеме́нт}. In the \ru{сложе́ние} class the neuter suffix \ru{-е́ние} can be added to a great variety of verbs to produce the corresponding abstract noun; thus \ru{дели́ть} \emph{to divide}, \ru{деле́ние} \emph{division}; \ru{написа́ть} \emph{to write}, \ru{написа́ние} \emph{writing}; \ru{относи́ть} \emph{to relate}, \ru{отноше́ние} \emph{relation} (note the consonant-variation in \ru{-нос-, -нош-}); \ru{производи́ть} \emph{produce}, \ru{произведе́ние} \emph{product} (note the vowel-gradation in \ru{-вод-, -вед-}); \ru{сумми́ровать} \emph{to sum}, \ru{сумми́рование} \emph{summation}, etc.

The types of nouns of next greatest frequency are the ones represented by \ru{коммутати́вность} and \ru{су́мма}. Most of the nouns in the \ru{коммутати́вность} class end either in \ru{-ность} \emph{-ity}: e.g. \ru{ассоциати́вность} \emph{associativity}, \ru{транзити́вность} \emph{transitivity}, or in \ru{-мость} \emph{-ability}: e.g. \ru{сумми́руемость} \emph{summability}. The class to which a noun belongs can always be recognized from its dictionary-form (i.e. its nominative singular) except for nouns ending in \ru{-ь}, which may belong either to the (in mathematics rather rare) \ru{числи́тель} class (masculine) or to the \ru{коммутати́вность} class (feminine).

The \ru{сумма} class contains several native Russian words, e.g. \ru{то́чка} \emph{point}, etc., and many transliterations, e.g. \ru{асимпто́та, директри́са, координа́та, систе́ма}.

Only one of the other classes, the \ru{место} class, contains any considerable number of words common in mathematics, most of them ending in the abstract suffix \ru{-ство}, e.g. \ru{доказа́тельство} \emph{proof}.

The \ru{числи́тель} class consists chiefly of words with the suffix \ru{-тель} \emph{-tor}; e.g. \ru{знамена́тель} \emph{denominator}.

The \ru{фу́нкция} class consists almost entirely of words ending in the suffix \ru{-ция} \emph{-tion}; e.g. \ru{инду́кция, констру́кция, опера́ция, прое́кция}.

In pure mathematics, the word \ru{по́ле} \emph{field} is the only representative of its class.

\section{Declension in the plural}

The endings for the plural of nouns are

\begin{center}
\small
\begin{tabular}{@{}lcc@{}}
\toprule
& hard/back patterns & soft/front patterns\\
\midrule
nom. & \ru{-ы/-а} & \ru{-и/-я}\\
gen. & \ru{-ей/-ов/-ий/---} & \\
dat. & \ru{-ам} & \ru{-ям}\\
acc. & \ru{-ы/-а} & \ru{-и/-я}\\
instr. & \ru{-ами} & \ru{-ями}\\
loc. & \ru{-ах} & \ru{-ях}\\
\bottomrule
\end{tabular}
\end{center}

where the symbol \ru{-ы/-а} means that masculine and feminine nouns end in \ru{-ы} and neuters in \ru{-а}. Here again the spelling-rule demands \ru{и} instead of \ru{ы} after \ru{к, г, х} or a sibilant; e.g. \ru{кни́ги} \emph{books}.

It was stated above that in the plural the three declensions have coalesced almost completely. In fact, earlier Russian had at least four types of declension in the plural, as is shown by the four different endings (front, back, front-reduced, back-reduced) for the genitive plural: \ru{-ей} in \ru{числи́телей}, \ru{-ов} in \ru{фа́ктов}, \ru{-ий} in \ru{сложе́ний} and --- (i.e. nothing at all) in \ru{сумм} and \ru{мест}. Note that this zero ending was represented in the old spelling by the hard sign \ru{ъ}, which was dropped from the writing in 1918 (see Chap. I, §4), and that some of these ``curtailed'' genitive plurals insert a movable \ru{е} or \ru{о} (see above), e.g. \ru{чи́сел} gen. plur. of \ru{число́} \emph{number}.

The declension in the plural of the eight types of nouns is as follows (to be practiced aloud):

\begingroup
\footnotesize
\setlength{\tabcolsep}{3pt}
\renewcommand{\arraystretch}{1.05}
\begin{longtable}{@{}l >{\cyrillicfont}p{23mm} >{\cyrillicfont}p{23mm} >{\cyrillicfont}p{25mm}@{}}
\toprule
& \normalfont facts & \normalfont loci & \normalfont sums\\
\midrule
nom. & факты & места & суммы\\
gen. & фактов & мест & сумм\\
dat. & фактам & местам & суммам\\
acc. & факты & места & суммы\\
instr. & фактами & местами & суммами\\
loc. & (о) фактах & (о) местах & (о) суммах\\
\midrule
& \normalfont numerators & \normalfont fields & \normalfont additions\\
\midrule
nom. & числители & поля & сложения\\
gen. & числителей & полей & сложений\\
dat. & числителям & полям & сложениям\\
acc. & числители & поля & сложения\\
instr. & числителями & полями & сложениями\\
loc. & (о) числителях & (о) полях & (о) сложениях\\
\midrule
& \normalfont functions & \multicolumn{2}{c}{commutativities}\\
\midrule
nom. & функции & \multicolumn{2}{>{\cyrillicfont}l}{коммутативности}\\
gen. & функций & \multicolumn{2}{>{\cyrillicfont}l}{коммутативностей}\\
dat. & функциям & \multicolumn{2}{>{\cyrillicfont}l}{коммутативностям}\\
acc. & функции & \multicolumn{2}{>{\cyrillicfont}l}{коммутативности}\\
instr. & функциями & \multicolumn{2}{>{\cyrillicfont}l}{коммутативностями}\\
loc. & (о) функциях & \multicolumn{2}{>{\cyrillicfont}l}{(о) коммутативностях}\\
\bottomrule
\end{longtable}
\endgroup

\section{Remarks on the exercises}

We are now ready for the first of the numerous exercises that consist of mathematical sentences from the Reading Selections and similar material, with word-for-word translations. Whenever the student encounters a strange word, as will happen often at first, he should try, as was noted in the Introduction, to observe it carefully enough so that when it turns up in its proper setting in Chapter V, and again in the Readings, it will not seem totally unfamiliar. If he is already interested in how it is constructed (usually from a compound verb) he may look it up in the general Glossary, where he will find a reference to the systematic discussion of it in the chapter on vocabulary (Chapter V).

These exercises illustrate the truly remarkable similarity of Russian word-order to English (slightly distorted from time to time in the translations given below), which makes it unnecessary to include a special chapter on syntax. The few remarks that have been necessary, e.g. on the genitive of comparison, on conditional sentences, etc., can be found when desired by consulting the English Index at the end of the book.

\begin{keybox}[Exercise 3.1]
Read aloud the following sentences, or parts of sentences, and account for the case of each noun. It is particularly valuable to read the sentence aloud once, looking at the Russian words, and then again without looking at them, as a kind of ``momentary memorization.''
\end{keybox}

\begin{enumerate}[label=\arabic*)]
\item \ru{то́чка M име́ет координа́ты x, y}\\
      the point M has coordinates $x,y$.
\item \ru{отноше́ние равно́ числу́}\\
      the ratio is equal to the number.
\item \ru{рассма́триваем систе́му координа́т}\\
      we consider a system of coordinates.
\item \ru{положе́ние то́чки M определя́ется отре́зком}\\
      the position of the point M is determined by the segment.
\item \ru{ве́кторы за́даны прое́кциями}\\
      the vectors are given by their projections.
\end{enumerate}

\section{The concept of grammatical gender}

In the preceding sections we have tacitly implied that all nouns declined like \ru{факт} and \ru{числи́тель} are masculine, those like \ru{место, по́ле} and \ru{сложе́ние} are neuter, and those like \ru{су́мма, фу́нкция} and \ru{коммутати́вность} are feminine, without discussing the significance of this somewhat startling terminology. In fact, the so-called ``gender'' of a noun is determined not by its manner of declension (since the above rule has a few exceptions of no mathematical interest) but by the behavior of modifying adjectives. A Russian schoolteacher, wishing to remind a pupil that \ru{факт} is masculine, \ru{су́мма} feminine and \ru{место} neuter, will not ordinarily use these technical, sex-related terms, but will say something like ``it's \ru{тот факт, та су́мма} and \ru{то место}'' (that fact, that sum, that place) with emphasis on the varying form of \ru{тот, та, то} for the word \emph{that}. As was pointed out by the Greek grammarians, a noun is called \emph{masculine} solely because, in its effect on the ending of an accompanying adjective or pronoun, it happens to act in the same way as the names of most of the male animals. In reading mathematics the student should try to remain sensitive to this agreement in gender for the endings of nouns, pronouns and adjectives (especially participles), since it reflects logical interrelations. The sections on pronouns and adjectives will provide many examples.

\section{Declension of pronouns}

The following pronouns, declined below, are common in mathematics:

\begin{center}
\small
\begin{tabular}{@{}lll@{}}
reflexive & \ru{себя́} & ourselves, itself, themselves\\
personal, 1st person & \ru{мы} & we\\
personal, 3rd person & \ru{он, они́} & it, they\\
demonstrative & \ru{тот, э́тот} & that, this\\
interrogative & \ru{что} & what\\
\end{tabular}
\end{center}

Pronouns in the Slavic languages fall into two groups: those that have remained without influence on the declension of adjectives, and those whose endings have been transferred to adjectives. The two pronouns \ru{себя́} and \ru{мы}, with distinctive declensions of their own, belong to the first class, and the others \ru{он, тот, этот, что}, all of them declined very much alike, belong to the second.

\begin{center}
\small
\begin{tabular}{@{}l >{\cyrillicfont}l@{}}
\multicolumn{2}{c}{\textit{Reflexive pronoun: singular and plural}}\\
nom. & ---\\
gen. & себя\\
dat. & себе\\
acc. & себя\\
instr. & собой\\
loc. & (о) себе\\
\end{tabular}
\end{center}

Note that the accusative of the reflexive (which obviously has no nominative) is like the genitive, whereas for nouns it was like the nominative. This identity of accusative and genitive, which for pronouns goes back to Proto-Slavic, was gradually extended, from about 1300 to 1800 A.D., to names of male animate beings. For example, the accusative singular of \ru{Ги́льберт} (Hilbert), otherwise declined like \ru{факт}, is \ru{Ги́льберта}.


\begin{center}
\small
\begin{tabular}{@{}l >{\cyrillicfont}l@{}}
\multicolumn{2}{c}{\textit{Declension of \ru{мы} ``we''}}\\
nom. & мы\\
gen. & нас\\
dat. & нам\\
acc. & нас\\
instr. & нами\\
loc. & (о) нас\\
\end{tabular}
\end{center}

Here it is clear that two pronouns have been combined, \ru{мы} for the nominative case, and elsewhere some other pronoun, whose nominative has been lost. The same phenomenon, common also in English (cf. \emph{we, us; it, they} etc.), occurs again in the third person pronoun \ru{он} etc., declined as follows.

The third person pronoun \ru{он} etc. is declined as follows:

\begin{center}
\small
\begin{tabular}{@{}l >{\cyrillicfont}l >{\cyrillicfont}l >{\cyrillicfont}l >{\cyrillicfont}l@{}}
\toprule
& masculine & neuter & feminine & plural\\
\midrule
nom. & он & оно & она & они\\
gen. & его & его & её & их\\
dat. & ему & ему & ей & им\\
acc. & его & его & её & их\\
instr. & им & им & ей & ими\\
loc. & (о) нём & (о) нём & (о) ней & (о) них\\
\bottomrule
\end{tabular}
\end{center}

Note that here, and throughout the book, we translate \ru{он} (masc.) and \ru{она́} (fem.) by \emph{it}, on the assumption that in mathematical passages the reference will be to a masculine or feminine noun like \ru{факт} \emph{fact} or \ru{су́мма} \emph{sum}. In the relatively rare instances (i.e. rare in mathematics) where \ru{он} and \ru{она́} refer to a person they must, of course, be translated by \emph{he} and \emph{she}. Similarly, the 3rd person reflexive pronoun \ru{себя́} etc. can mean, not only \emph{itself} or \emph{themselves}, but also \emph{himself, herself} or \emph{ourselves}.

Note also that when these pronouns are governed by a preposition (e.g. \ru{у, к, с, в, о}, etc.) the letter \ru{н} is prefixed to them, as shown above for the locative; thus \ru{у него́} \emph{at it}, \ru{к нему́} \emph{toward it}, \ru{в нём} \emph{in it}; for brevity, we omit the historical explanation.

The nominatives \ru{он, она́, оно́, они́} have had no influence elsewhere in Russian declension, but the endings of the other cases \ru{-его, -ему, ...} (\ru{-ого, -ому}, etc. after a hard consonant) occur in all the other pronouns and in almost all adjectives. In modern Russian the \ru{-г-} in the genitive singular of such pronouns and adjectives is pronounced like \ru{в}; for example, \ru{чего́} \emph{of what} is pronounced approximately \emph{chee-VO}, the spelling with \ru{г} being due to the influence of Church Slavonic, where such forms were both pronounced and written with \ru{г}.\corrnote{原文在这里进一步声称这种 \ru{-ого/-его} 的 [v] 读音“在文字发明前的古俄语中就已经存在”。这一年代判断并不可靠；故仅删去该错误的年代说明，其余原文论述保留。}

Thus we can set up the following tables.

\begin{center}
\small
\begin{tabular}{@{}l >{\cyrillicfont}l l@{}}
\multicolumn{3}{c}{\textit{Interrogative pronoun}}\\
nom. & что & what\\
gen. & чего & \\
dat. & чему & \\
acc. & что & \\
instr. & чем & \\
loc. & (о) чём & \\
\end{tabular}
\end{center}

The negative pronoun \ru{ничто́} \emph{nothing} is declined like \ru{что}:

\begin{center}
\small
\begin{tabular}{@{}l >{\cyrillicfont}l l >{\cyrillicfont}l l@{}}
nom.   & ничто   & nothing    & ни для чего & for nothing\\
gen.   & ничего  & of nothing & ни к чему   & toward nothing\\
dat.   & ничему  & to nothing &             &\\
acc.   & ничто   &            &             &\\
instr. & ничем   &            & ни с чем    & with nothing\\
loc.   &         &            & ни о чём    & about nothing\\
\end{tabular}
\end{center}

Note that when \ru{ничто́} is used with a preposition, the latter is placed between its two parts, so as to form a phrase of three separate words; i.e. \emph{about nothing} becomes \emph{not about anything}: \ru{ни о чём}.

The pronominal adjectives \ru{тот} \emph{that}, \ru{э́тот} \emph{this} and \ru{весь} \emph{all} are declined as follows. Note that \ru{э́тот} has \ru{э-} where \ru{тот} has \ru{о-}, and that the stem \ru{всь-} of \ru{весь} ends in a soft consonant, which will therefore be followed by \ru{-е-} rather than \ru{-о-}. The \ru{е} in the masculine nominative singular \ru{весь} is a movable vowel (see §2).

\begingroup
\footnotesize
\setlength{\tabcolsep}{3.2pt}
\begin{longtable}{@{}l >{\cyrillicfont}p{26mm} >{\cyrillicfont}p{26mm} >{\cyrillicfont}p{26mm}@{}}
\toprule
& \normalfont \ru{тот} & \normalfont \ru{этот} & \normalfont \ru{весь}\\
\midrule
\multicolumn{4}{@{}l}{\textit{Masculine/neuter}}\\
nom. & тот / то & этот / это & весь / всё\\
gen. & того & этого & всего\\
dat. & тому & этому & всему\\
acc. & тот / то & этот / это & весь / всё\\
instr. & тем & этим & всем\\
loc. & (о) том & (об) этом & (о) всём\\
\midrule
\multicolumn{4}{@{}l}{\textit{Feminine}}\\
nom. & та & эта & вся\\
gen. & той & этой & всей\\
dat. & той & этой & всей\\
acc. & ту & эту & всю\\
instr. & той & этой & всей\\
loc. & (о) той & (об) этой & (о) всей\\
\midrule
\multicolumn{4}{@{}l}{\textit{Plural}}\\
nom. & те & эти & все\\
gen. & тех & этих & всех\\
dat. & тем & этим & всем\\
acc. & те & эти & все\\
instr. & теми & этими & всеми\\
loc. & (о) тех & (об) этих & (о) всех\\
\bottomrule
\end{longtable}
\endgroup


The neuter nominative singular \ru{то} occurs commonly in the abbreviation \ru{т.е.} = \ru{то есть}, ``i.e., that is.'' Similar in declension to \ru{этот} is \ru{оди́н} ``one'' (the numeral):

\begin{center}
\small
\begin{tabular}{@{}l >{\cyrillicfont}l >{\cyrillicfont}l@{}}
& masculine/neuter & feminine\\
nom. & один / одно & одна\\
gen. & одного & одной\\
dat. & одному & одной\\
\end{tabular}
\end{center}

\begin{keybox}[Exercise 3.2]
Read aloud the following sentences, noting the inflectional endings of the pronouns and pronominal adjectives.
\end{keybox}

\begin{enumerate}[label=\arabic*)]
\item \ru{это и есть иско́мая фо́рмула}\\
      this is the desired formula.
\item \ru{е́сли то́чка A лежи́т на пара́боле, то её координа́ты \ldots}\\
      if the point A lies on the parabola, then its coordinates \ldots
\item \ru{ве́ктор, кото́рый име́ет длину́, ра́вную b \ldots}\\
      a vector which has a length equal to $b$ \ldots\\[-1mm]
      \emph{(Note that in Russian a mathematical symbol can have a definite case; here the symbol $b$ is in the dative case requiring the insertion of ``to'' in English.)}
\item \ru{она́ лежи́т на э́ллипсе в том и то́лько в том слу́чае, когда́ \ldots}\\
      it lies on the ellipse in that and only in that case, when \ldots\\[-1mm]
      \emph{(Note this expression for ``if and only if.'')}
\item \ru{мы ви́дим, что \ldots}\\
      we see that \ldots\\[-1mm]
      \emph{(Recall that \ru{что} is pronounced approximately \ru{што}.)}
\item \ru{существу́ет така́я окре́стность э́той то́чки, в кото́рой \ldots}\\
      there exists such a neighborhood of this point, in which \ldots\\[-1mm]
      \emph{(Note the word order for ``a neighborhood such that \ldots'')}
\end{enumerate}

\section{Declension of adjectives}

The endings for the adjectives were taken over from the pronouns, as stated above. In many instances they represent a fusion of noun and pronoun endings.

\begin{center}
\small
\begin{tabular}{@{}l >{\cyrillicfont}l >{\cyrillicfont}l >{\cyrillicfont}l@{}}
\toprule
& masculine/neuter & feminine & plural, all genders\\
\midrule
nom. & -ый / -ое & -ая & -ые\\
gen. & -ого & -ой & -ых\\
dat. & -ому & -ой & -ым\\
acc. & -ый / -ое & -ую & -ые\\
instr. & -ым & -ой & -ыми\\
loc. & -ом & -ой & -ых\\
\bottomrule
\end{tabular}
\end{center}

For example, the adjective \ru{ра́вный} \emph{equal} is declined:

\begin{center}
\small
\begin{tabular}{@{}l >{\cyrillicfont}l >{\cyrillicfont}l >{\cyrillicfont}l@{}}
\toprule
& masculine/neuter & feminine & plural\\
\midrule
nom. & равный / равное & равная & равные\\
gen. & равного & равной & равных\\
dat. & равному & равной & равным\\
acc. & равный / равное & равную & равные\\
instr. & равным & равной & равными\\
loc. & равном & равной & равных\\
\bottomrule
\end{tabular}
\end{center}

Let us here make six remarks.

\begin{enumerate}[label=\roman*)]
\item If the ending is accented, the \ru{-ый} of the masculine nominative singular becomes \ru{-о́й}, e.g. \ru{второ́й} \emph{second}, \ru{просто́й} \emph{simple}, \ru{пусто́й} \emph{empty}, \ru{непусто́й} \emph{non-empty}, etc.

\item After \ru{к, г, х} or a sibilant, the spelling-rule again demands \ru{и} instead of \ru{ы}; e.g. \ru{коро́ткий} \emph{short}, \ru{коро́ткая, коро́ткие}, etc.; \ru{вся́кий} \emph{every}, etc. Note also \ru{больши́м} (not *\ru{большы́м}) in the instrumental of \ru{большо́й} \emph{big}, and similarly in \ru{небольшо́й} \emph{not big}. (Compare also the noun \ru{большинство́} \emph{majority}.)

\item Just as in the English sentence \emph{the good die young}, so in Russian many adjectives can be used as nouns; e.g. the word \ru{произво́дная} \emph{derivative} is a feminine adjective modifying the noun \ru{фу́нкция} \emph{function} understood.

\item When an adjective is used predicatively, it may appear in a short form obtained by dropping the ending of the masculine singular, e.g. \ru{прост, проста́, про́сто} sing.; \ru{просты́} plur. Thus \ru{просто́й факт} \emph{a simple fact}, but \ru{тот факт прост} \emph{that fact is simple}. If the stem of the adjective ends in two consonants considered difficult to pronounce together, a movable \ru{е} or \ru{о} is inserted in the short form of the nom. masc. sing.; thus \ru{ра́вен, равна́, равно́; равны́} \emph{equal}.

For an adjective formed from a proper name the short forms may also be used attributively, e.g. \ru{евкли́дова геоме́трия} \emph{Euclidean geometry}.

\item In the declension of nouns, we have seen that stems ending in a hard consonant (i.e. of the types \ru{факт, место, су́мма}) are much more numerous than their counterparts with a soft consonant (\ru{числи́тель, по́ле, фу́нкция}). With adjectives the preference for hard endings is even more noticeable.

\begin{sloppypar}
In fact, the one hard-consonant suffix \ru{-н-} accounts for well over three-quarters of all adjectives used in mathematics. Thus \ru{аналогичный} \emph{analogous}, \ru{ве́рный} \emph{true}, \ru{действи́тельный} \emph{real}, \ru{доста́точный} \emph{sufficient}, \ru{еди́нственный} \emph{unique}, \ru{положи́тельный} \emph{positive}, \ru{поля́рный} \emph{polar}, \ru{пропорциона́льный} \emph{proportional}, \ru{симметри́чный} \emph{symmetric}, \ru{фундамента́льный} \emph{fundamental}, \ru{эквивале́нтный} \emph{equivalent}, \ru{элемента́рный} \emph{elementary}.
\end{sloppypar}

Most other adjectives have a stem ending in some other hard consonant (and are therefore declined like \ru{ра́вный}), e.g. \ru{необходи́мый} \emph{necessary}, \ru{но́вый} \emph{new}, \ru{одина́ковый} \emph{identical}, \ru{пе́рвый} \emph{first}, etc. To this class also belong the various pronominal adjectives: i.e. the relative \ru{кото́рый} \emph{which}, the interrogative \ru{како́й} \emph{of what kind}, the demonstrative \ru{тако́й} \emph{of such a kind}, the personal \ru{наш} \emph{our}, the reflexive \ru{свой} \emph{one's own} (with the derived noun \ru{сво́йство} \emph{property} and adjective \ru{со́бственный} \emph{one's own, proper, eigen-}, etc.) and the indefinites \ru{ка́ждый} \emph{each}, \ru{не́который} \emph{some}, and \ru{никако́й} \emph{of any kind at all} (only in negative sentences). In prepositional phrases the \ru{ни-} in \ru{никако́й} becomes separated from the \ru{како́й} (cf. the note on \ru{ни о чём} in §7), thus: \ru{ни в како́м по́ле не вло́жится} \emph{it is not embedded in any field}; \ru{нигде́ не пло́тный} \emph{nowhere dense} (lit. anywhere not-dense; \ru{где} \emph{where}, \ru{пло́тный} \emph{dense}).\corrnote{原例印作 \ru{нигде неплотный}；数学术语通常写作 \ru{нигде не плотный}，故作最小正字法修正。}

Such adjectives (or pronominal adjectives) are declined as follows:

\begin{center}
\footnotesize
\begin{tabular}{@{}>{\cyrillicfont}p{40mm} >{\cyrillicfont}p{40mm} >{\cyrillicfont}p{40mm}@{}}
\normalfont masculine/neuter & \normalfont feminine & \normalfont plural\\
который / которое & которая & которые\\
которого & которой & \ldots\\
\ldots & \ldots & \ldots\\[1mm]
такой / такое & такая & такие\\
такого & такой & \ldots\\
\ldots & \ldots & \ldots\\[1mm]
наш / наше & наша & наши\\
нашего & нашей & \ldots\\
\ldots & \ldots & \ldots\\
\end{tabular}
\end{center}

Note that some words have a prefix \ru{не-} which is of different origin and meaning from the negative prefix; thus \ru{не́который} means \emph{some}. Note also how the spelling-rule demands \ru{и} instead of \ru{ы} in the plural forms \ru{таки́е} \emph{such} and \ru{на́ши} \emph{our}.

\item But in spite of the above preference for hard endings, there are a few adjectives with stems ending in soft \ru{н} (and one or two, e.g. \ru{о́бщий} \emph{general}, in some other soft consonant), which means, of course, that the endings must be \ru{-ий, -его, -ему}, etc. Adjectives of this type are \ru{ве́рхний} \emph{upper}, \ru{вну́тренний} \emph{interior}, \ru{ни́жний} \emph{lower}, \ru{после́дний} \emph{latter}, \ru{сре́дний} \emph{average}, \ru{тре́тий} \emph{third}. For example,
\end{enumerate}

\begin{center}
\small
\begin{tabular}{@{}l >{\cyrillicfont}l >{\cyrillicfont}l@{}}
\toprule
& masculine/neuter & feminine\\
\midrule
nom. & последний / последнее & последняя\\
gen. & последнего & последней\\
dat. & последнему & последней\\
acc. & последний / последнее & последнюю\\
instr. & последним & последней\\
loc. & последнем & последней\\
\bottomrule
\end{tabular}
\end{center}

\begin{keybox}[Exercise 3.3]
Read the following sentences aloud, noting the various adjectives, especially those in the short predicative form.
\end{keybox}

\begin{enumerate}[label=\arabic*)]
\item \ru{таки́е ве́кторы равны́}\\
      such vectors (are) equal.
\item \ru{пусть M --- произво́льная то́чка}\\
      let M (be) an arbitrary point.
\item \ru{так как ве́кторы паралле́льны, то их координа́ты пропорциона́льны}\\
      since the vectors are parallel, therefore their coordinates are proportional.
\item \ru{мы изло́жим тео́рию иррациона́льных чи́сел}\\
      we shall explain the theory of irrational numbers.
\item \ru{одна́ко э́тот необходи́мый при́знак \ldots}\\
      but this necessary criterion \ldots
\item \ru{интегра́льное исчисле́ние ста́вит пе́ред собо́й обра́тную зада́чу}\\
      the integral calculus sets before itself the inverse problem.
\end{enumerate}

\begin{keybox}[Exercise 3.4]
Read the following sentences aloud, noting the various pronominal forms.
\end{keybox}

\begin{enumerate}[label=\arabic*)]
\item \ru{систе́ма координа́т, нача́ло кото́рой совпада́ет с по́люсом}\\
      a system of coordinates, the origin of which coincides with the pole.
\item \ru{ве́ктор, кото́рый определя́ется усло́виями}\\
      a vector which is defined by the conditions.
\item \ru{мо́жно счита́ть их ко́синусом и си́нусом не́которого угла́}\\
      it is possible to consider them as the cosine and sine of some angle.
\item \ru{существу́ет така́я то́чка, что \ldots}\\
      there exists such a point that \ldots
\item \ru{G не мо́жет быть о́бластью голомо́рфности никако́й фу́нкции}\\
      G cannot be the domain of holomorphy of any function.
\item \ru{тепе́рь мы ви́дим, каку́ю пробле́му изуча́ем}\\
      now we see what sort of problem we are studying.
\end{enumerate}


\section{The numerals}

In Russian, as in English, the ordinal numerals are declined like adjectives; thus \ru{пе́рвый} \emph{first}, \ru{второ́й} \emph{second}, \ru{тре́тий} \emph{third}.

But the cardinals behave in such an unusual way that the Russian grammarians even call the numerals a separate part of speech. In fact, apart from \ru{оди́н} \emph{one} (see §7) they behave more like nouns than like adjectives. Thus \ru{два} \emph{two} is declined as follows:

\begin{center}
\small
\begin{tabular}{@{}l >{\cyrillicfont}l l@{}}
\toprule
nom. & два (fem. две) & two\\
gen. & двух & of two\\
dat. & двум & to two\\
acc. & два (fem. две) & two\\
instr. & двумя & by/with two\\
loc. & (о) двух & about two\\
\bottomrule
\end{tabular}
\end{center}

and the form of the noun after the word \emph{two} looks like a genitive singular (though in fact it is historically connected with the old Indo-European dual, used for a pair of things), e.g. \ru{два квадра́та} \emph{two squares}.\corrnote{原文直接写作 “though in fact it is an old Indo-European dual”。现代俄语中该形式同步上不是一个仍有生产力的双数范畴；这里作最小修正为“historically connected with the old dual”。} However, a modifying adjective is put in the plural; e.g. \ru{два больши́х квадра́та} \emph{two large squares}, and this somewhat startling usage is extended to the cardinals \ru{три, трёх, ...} \emph{three} and \ru{четы́ре, четы́рёх, ...} \emph{four}. But beginning with \ru{пять} \emph{five} (and in some modern writers even beginning with \ru{два} \emph{two}) the noun is also in the genitive plural; thus \ru{пять квадра́тов} \emph{five squares} (more literally, a quintet of squares). Fortunately, in a mathematical context the meaning is always clear.

\section{Adjective--noun phrases; italics}

The following exercise, consisting of adjective-noun phrases (e.g. \ru{математи́ческая ло́гика} \emph{mathematical logic}) taken from the Subject-Classification of the Russian counterpart of \emph{Mathematical Reviews}, namely \ru{Реферати́вный журна́л --- Матема́тика}, will illustrate how an adjective agrees in gender, number and case with the noun it modifies. The phrases are given here in italics, in order to provide the practice necessary for reading the statements of theorems etc. in an ordinary Russian mathematical text.

Only six of the italic letters are enough different from their non-italicized counterparts to cause any trouble; namely italic \ru{г} (No. 4), \ru{д} (No. 5, but a roman-like form is also used), \ru{и} (No. 10), \ru{й} (No. 11), \ru{п} (No. 17), and \ru{т} (No. 20, but a roman-like \ru{т} is sometimes used). The Russian italic letters originated from Latin handwriting; e.g. the italic form of \ru{т} is the final result of a long series of efforts to write \ru{т} unicursally, i.e. without lifting the pen (the italic form of \ru{м} is also notably different).

The mathematical terms in the exercise will all be familiar, since they are transliterations of international words, but the reader should look carefully at the agreement (in gender, number and case) of the adjective with the noun.

\begin{keybox}[Exercise 3.5]
Read aloud the following adjective--noun phrases, printed in italics.
\end{keybox}

\begingroup
\cyrillicfont\itshape\small\sloppy
Математическая логика, дедуктивные системы, модальная логика,
алгебраическая логика, комбинаторная логика.

Комбинаторный анализ, биномиальные коэффициенты, комбинаторные
функции, ортогональные таблицы, латинские квадраты.

Коммутативные алгебры, гомологические методы, алгебраическая
геометрия, абстрактные дифференцирования, позитивные матрицы,
квадратичные формы, мультилинейная алгебра, ассоциативные алгебры,
гомологическая алгебра, абстрактные категории, гомологическая алгебра.

Периодические группы, гомологические методы, классические группы,
бинарные системы, топологические группы.

Элементарные функции, монотонные функции, алгебраические функции,
модулярные функции, экстремальные проблемы, супергармонические функции,
специальные функции, эллиптические функции, гипергеометрические
функции, цилиндрические функции, вариационные методы,
псевдодифференциальные операторы, квазилинейные системы.

Тригонометрическое интерполирование, ортогональные функции,
тригонометрические проблемы, спектральный синтез, функциональный анализ,
спектральная теория, компактные операторы, интегральные операторы.

Специальные геометрии, дифференциальная геометрия, тензорный анализ,
дифференциальные инварианты, локальная теория, геометрические объекты,
глобальная теория, интегральная геометрия.

Локальная компактность, топологическая динамика, алгебраическая
топология, когомологические операции, геометрическая топология,
дифференцируемые структуры, характеристические классы.

Стохастические процессы, математическая статистика, асимптотическая
теория, оптимальные программы, непараметрические методы.

Математическое программирование, графические методы, гармонический
анализ, классическая термодинамика, квантовая механика,
статистическая физика, математическая экономика, динамическое
программирование.
\par
\endgroup

\section{Comparative and superlative of adjectives and adverbs}

The comparative degree of an adjective is formed, as in English, in three different ways:
\begin{enumerate}[label=\roman*)]
\item by prefixing the (uninflected) word \ru{бо́лее} \emph{more}: e.g. \ru{бо́лее элемента́рный} \emph{more elementary};
\item by adding the (uninflected) suffix \ru{-ее} (cf. the English \emph{-er}): e.g. \ru{старе́е} \emph{older} (from \ru{ста́рый} \emph{old}), which in words involving ``consonant variation'' (see §6 of the Introduction) becomes \ru{-е}, e.g. \ru{про́ще} \emph{simpler} from \ru{просто́й} \emph{simple}; \ru{бо́льше} \emph{bigger} from \ru{большо́й} \emph{big}, etc. (cf. the dropping of one of the e's with the English word \emph{simpler});
\item for a few adjectives (as in the English \emph{good, better, best}), the comparative (and superlative) degrees do not come from the same root as the positive: e.g. \ru{хоро́ший} \emph{good}, \ru{лу́чше} \emph{better}; \ru{ма́лый} \emph{small}, \ru{ме́ньше} (or \ru{ме́нее}) \emph{less}, \ru{ме́ньший} \emph{least}.\corrnote{原文把 \ru{лучше} 注为 “better (or best)”。现代标准俄语中 \ru{лучше} 是比较级“better”，而“best”通常用 \ru{лу́чший} 或分析式表达，故作最小更正。}
\end{enumerate}

The superlative degree may be formed in two different ways (again compare the English):
\begin{enumerate}[label=\roman*)]
\item by prefixing the (inflected) adjective \ru{са́мый} \emph{most, very}; e.g. \ru{са́мый элемента́рный} \emph{most elementary};
\item by adding the (inflected) suffix \ru{-е́йший} (English \emph{-est}), e.g. \ru{нове́йший} \emph{newest} (from \ru{но́вый} \emph{new}), \ru{простейший} \emph{simplest} (from \ru{просто́й} \emph{simple}), which in words involving consonant variation (Introduction §6) becomes \ru{-а́йший}, e.g. \ru{широ́кий} \emph{wide}, \ru{широча́йший} \emph{widest}; and which is sometimes shortened, e.g. \ru{ста́рый} \emph{old}, \ru{ста́рший} \emph{oldest}.
\end{enumerate}

The superlative can be strengthened by prefixing \ru{наи-}; e.g. \ru{наибо́льший} \emph{the very largest}, \ru{наиме́ньший} \emph{the very least}.

Certain adverbs, namely the so-called adverbs of quality formed directly from adjectives (in English, mostly by adding \emph{-ly}) are compared like the corresponding adjectives. For most of them the positive degree is like the accusative neuter singular of the short (predicative) form of the adjective (§8, Remark iv); e.g. \ru{аналоги́чный} \emph{analogous}, \ru{аналоги́чно} \emph{analogously}; but if the adjective ends in \ru{-ский} the adverb ends in \ru{-ски}; thus \ru{математи́ческий} \emph{mathematical}, \ru{математи́чески} \emph{mathematically}. The corresponding comparatives and superlatives are rarely used in mathematics, e.g. \ru{бо́лее про́сто} \emph{more simply}.

\begin{keybox}[Exercise 3.6]
Read aloud the following sentences, noting the comparative and superlative forms.
\end{keybox}

\begin{enumerate}[label=\arabic*)]
\item \ru{э́то число́ бу́дет ли́бо наибо́льшим в ни́жнем кла́ссе, ли́бо наиме́ньшим в ве́рхнем кла́ссе}\
      this number will be either the largest in the lower class or the smallest in the upper class.
\item \ru{простейшими из логи́ческих исчисле́ний явля́ются исчисле́ния выска́зываний, класси́ческое и конструкти́вное}\
      the simplest of the logical calculi are the calculi of propositions, classical and constructive.
\item \ru{счита́ют бо́лее вероя́тным, что \ldots}\
      they consider it more probable that \ldots
\item \ru{зада́ча сво́дится к тако́му вы́бору фу́нкции, при кото́ром интегра́л име́ет наиме́ньшее значе́ние}\
      the problem reduces to such a choice of the function for which the integral has the least value.
\end{enumerate}

\section{The uninflected parts of speech}

Apart from the interjection (of no importance in mathematics) the Russian grammarians list four uninflected parts of speech: i) adverb, ii) conjunction, iii) particle, and iv) preposition.

\subsection*{Adverbs}
The chief adverbs (excepting those formed directly from adjectives) are
\begin{center}
\small
\begin{tabular}{@{}>{\cyrillicfont}l l >{\cyrillicfont}l l >{\cyrillicfont}l l@{}}
тепе́рь & now & тогда́ & then & когда́ & when\\
здесь & here & там & there & где & where\\
сюда́ & hither & туда́ & thither & куда́ & whither\\
отсю́да & hence & отту́да & thence & отку́да & whence\\
всю́ду & everywhere & так & thus, so & как & how\\
ина́че & otherwise & иногда́ & sometimes & всегда́ & always\\
& & уже́ & already & &
\end{tabular}
\end{center}
and some single words that were originally phrases, like \ru{наприме́р} \emph{for example} (\ru{приме́р} \emph{example}).

\subsection*{Conjunctions}
The chief subordinating conjunctions, preceded in Russian by a comma, are \ru{что} \emph{that}, \ru{чем} \emph{than}, \ru{е́сли} \emph{if}, \ru{так как} \emph{since, because}, \ru{что́бы} \emph{in order that}, \ru{ввиду́ того́, что} \emph{in view of the fact that}, \ru{потому́ что} \emph{because} (lit. in accordance with the fact that), \ru{причём} \emph{where} (lit. in the presence of which).

The chief coordinating conjunctions are \ru{и} \emph{and}, \ru{и ... и} (also \ru{так ... как}) \emph{both ... and}, \ru{ита́к} \emph{and so}, \ru{одна́ко} \emph{nevertheless}, \ru{но} \emph{but}, \ru{а} \emph{and or but}, \ru{и́ли} \emph{or}, \ru{ли́бо ... ли́бо} \emph{either ... or}, \ru{ни ... ни} \emph{neither ... nor}, \ru{поэ́тому} \emph{therefore} (lit. in accordance with this), \ru{прито́м} \emph{moreover} (lit. in the presence of that; also \ru{при э́том} as two words), \ru{то} \emph{therefore}.

\subsection*{Particles}
Chief among the words called particles are \ru{не} \emph{not}; \ru{то} \emph{that}, as in \ru{то есть} \emph{that is}, abbreviated \ru{т.е.}; \ru{же}, which emphasizes the preceding word; \ru{и}, in a special emphatic use, which emphasizes the following word; \ru{ли} \emph{whether} in questions; \ru{то́лько} \emph{only}; \ru{бы} \emph{would}; and \ru{пусть} \emph{let, let it be the case that, suppose that}. Note that \ru{пусть} is followed, not by an infinitive like the English \emph{let}, but by an independent clause; thus \emph{let it have} is \ru{пусть име́ет}, lit. \emph{let it has} (i.e. \emph{let it be the case that it has}).

Examples are: \ru{тот же детермина́нт} \emph{the same determinant} (note the derived words \ru{то́ждество} \emph{identity}, \ru{тожде́ственный} \emph{identical}, \ru{тожде́ственно} \emph{identically}, etc.); \ru{тогда́ же} \emph{at the same time}, \ru{так же} \emph{in the same way} (written as one word \ru{та́кже} means \emph{also}); \ru{е́сли же} \emph{but if}.

Further examples are:
\begin{quote}
\ru{э́то и есть иско́мая фо́рмула} --- this is in fact the desired formula;\\
\ru{как узна́ть, рациона́льна ли да́нная крива́я} --- how to recognize whether a given curve is rational (note the word-order);\\
\ru{тогда́ и то́лько тогда́, когда́ ...} --- then and only then, when ... (i.e. if and only if);\\
\ru{е́сли бы э́то бы́ло ве́рно, мы име́ли бы ...} --- if this were true, we would have ...
\end{quote}

\subsection*{Prepositions and prefixes}
Finally, the name \emph{preposition} (a Latin translation of the Greek word used by Aristotle) obviously refers to something \emph{placed before} something else. In English it is naturally taken to mean placed before a noun (as a separate word); but in Greek, Latin and Russian the meaning placed before a verb (as a prefix) is at least equally important. Those prepositions that can serve both functions are called \emph{proper}, those that can serve only the first are called \emph{improper}, and those that can serve only the second are not usually called prepositions at all, but \emph{inseparable prefixes}.

In Russian there are fifteen proper prepositions, not including \ru{к} \emph{toward}, which is very common before nouns and pronouns but has been so seldom prefixed to a verb that we shall neglect this use of it altogether (the forms in parentheses are used before words beginning with certain combinations of letters; e.g. \ru{обо мно́жестве} \emph{about the set}):
\begin{center}
\small
\begin{tabular}{@{}>{\cyrillicfont}l l >{\cyrillicfont}l l@{}}
до & to & за & after, in exchange for, as\\
из (изо) & from & на & on, onto\\
от (ото) & out of, from, of & под (подо) & under\\
с (со) & from & над (надо) & over\\
у & at & перед (передо) & in front of\\
по & according to & про & about\\
в (во) & in, into & о (об, обо) & about\\
при & in the presence of & &
\end{tabular}
\end{center}
and four inseparable prefixes: \ru{воз-/вз-} \emph{up}, \ru{вы-} \emph{out}, \ru{раз-} \emph{apart}, \ru{пере-} \emph{across}.

The improper prepositions, on the other hand, can hardly be listed in full, since new ones are constantly being formed. Among those that have been in the language for a relatively long time are \ru{без} \emph{without}, \ru{ввиду́} \emph{in view of}, \ru{вдоль} \emph{along}, \ru{вме́сте с} \emph{together with}, \ru{вме́сто} \emph{instead of}, \ru{вне} \emph{outside}, \ru{внутри́} \emph{inside}, \ru{всле́дствие} \emph{in consequence of}, \ru{для} \emph{for}, \ru{кро́ме} \emph{beyond, except}, \ru{ме́жду} \emph{between}, \ru{по́сле} \emph{after}, \ru{про́тив} \emph{against}, \ru{путём} \emph{by means of} (lit. by the path of), \ru{среди́} \emph{among}, \ru{че́рез} \emph{through, by means of}. Typical of the more recent ones is \ru{относи́тельно} \emph{relatively to, with respect to}. Almost all improper prepositions are followed by the genitive case; e.g. \ru{путём интегри́рования} \emph{by means of integration}, \ru{относи́тельно но́вой систе́мы координа́т} \emph{with respect to the new system of coordinates}, etc., but note e.g. \ru{ме́жду} followed by the instrumental: \ru{ме́жду то́чками} \emph{between the points}, and \ru{че́рез} followed by the accusative: \ru{пряма́я че́рез да́нную то́чку} \emph{a line through the given point}.

The proper prepositions are used with various cases:
\begin{center}
\small
\begin{tabular}{@{}l >{\cyrillicfont}p{112mm}@{}}
with genitive & до, из, от, с, у\\
with dative & к, по\\
with accusative & в, за, на, под, про\\
with instrumental & за, над, перед, под, с\\
with locative & в, на, о, при\\
\end{tabular}
\end{center}
\corrnote{原表的 locative 行只印有 \ru{в, на, при}，但同页前文明确列出 \ru{о/об/обо} ``about'' 且其支配前置格；此处补入 \ru{о}。}

As will be seen from this list and from the following exercises, each of these prepositions must be translated into English in various ways, depending on the context. Motion toward is usually expressed by a preposition with the accusative (sometimes with the dative), rest in with the locative (sometimes with the instrumental), and motion from with the genitive. Common to many Indo-European languages is the genitive of comparison: e.g. the line AB is longer from (i.e. than) the line CD, which is in fact an example of motion from; i.e. ``starting from (the end of) the line CD, the line AB is (by that amount) longer.'' For example, \ru{ка́ждое число́ мно́жества A ме́ньше ка́ждого числа́ мно́жества B}, \emph{each number of the set A is less than (lit. from) each number of the set B}; but the construction with \ru{чем} \emph{than}, as in English, is equally common: thus \ru{ме́ньше, чем ка́ждое число́} \emph{less than each number}.

\begin{keybox}[Exercise 3.7]
Read aloud these sentences illustrating prepositional phrases.
\end{keybox}

\begin{enumerate}[label=\arabic*)]
\item \ru{расстоя́ние от то́чки гипе́рболы до фо́куса}\\
      the distance from a point of the hyperbola to a focus.
\item \ru{наиме́ньшее из его́ дели́телей}\\
      the smallest of its divisors.
\item \ru{накло́н к о́си}\\
      inclination to the axis.
\item \ru{по определе́нию}\\
      by definition.
\item \ru{элеме́нт $b$ сле́дует за $a$}\\
      the element $b$ follows after $a$.
\item \ru{опера́ция над ве́кторами}\\
      operations on vectors.
\item \ru{ста́вим пе́ред собо́й зада́чу}\\
      we set before ourselves the problem.
\item \ru{нача́ло совпада́ет с по́люсом}\\
      the origin coincides with the pole.
\item \ru{деле́ние в за́данном отноше́нии}\\
      division in a preassigned ratio.
\item \ru{фу́нкция регуля́рна на криво́й, е́сли \ldots}\\
      a function is regular on a curve if \ldots
\item \ru{поня́тие о сече́нии}\\
      the concept of a cut.
\item \ru{преобразова́ние координа́т при паралле́льном перено́се}\\
      transformation of coordinates under parallel translation.
\item \ru{аналити́ческое продолже́ние вдоль криво́й}\\
      analytic continuation along a curve.
\item \ru{фу́нкция регуля́рна внутри́ о́бласти}\\
      the function is regular inside the domain.
\item \ru{здесь $\Delta y\to0$ вме́сте с $\Delta x$}\\
      here $\Delta y\to0$ together with $\Delta x$.
\item \ru{для ка́ждой то́чки име́ем \ldots}\\
      for each point we have \ldots
\item \ru{среди́ действи́тельных чи́сел нахо́дится число́, кото́рое \ldots}\\
      among the real numbers there is a number which \ldots
\item \ru{обозна́чим прое́кцию то́чки M на ось $Ox$ че́рез $M_x$}\\
      we shall denote the projection of the point M onto the axis $Ox$ by $M_x$.
\end{enumerate}

\section{The verb; present imperfective and future perfective}

In addition to tense, as in English, the regular verb in Russian has aspect, discussed just below and in the next chapter; and it has a complicated set of participles. But otherwise, at least for mathematics, it is simpler than in English: there are no special synthetic subjunctive or conditional forms (for the particle \ru{бы} in conditional and purpose clauses see Exercise 3.8 in §14); there is no separate synthetic passive conjugation, passive meanings being expressed chiefly by participles and reflexive forms; the imperative mood scarcely occurs in mathematics and the indicative occurs chiefly in the third person: \emph{it does, they do}, and in the first person plural: \emph{we do}; there are only three tenses, present, past and future, with no refinements like pluperfects or progressives; there are only two conjugations, represented by the infinitives \ru{чита́ть} \emph{to read} and \ru{дели́ть} \emph{to divide}, and a regular verb has two principal parts useful for the rules given here, namely the infinitive \ru{чита́ть, дели́ть} and the 1st person plural \ru{чита́ем} \emph{we read}, \ru{де́лим} \emph{we divide}.\corrnote{原文绝对写作“there are no subjunctive or conditional forms; there are no passive forms, except in the participles”。俄语用 \ru{бы}+过去式表达条件/虚拟意义，并可用 \ru{-ся} 结构表达被动；这里仅作最小纠正，同时保留作者针对数学语篇的教学框架。}

Most Russian verbs occur in pairs, one of the pair being an imperfective verb (or in the imperfective aspect) and the other a perfective verb (or in the perfective aspect). The difference in meaning between the two aspects will be discussed more fully in Chapter IV; here let us simply say that the imperfective aspect describes an action without reference to its completion, but the perfective implies something about its completion in the past or future, so that, in particular, a perfective verb cannot have a present tense.

For example, \ru{чита́ть} \emph{to read} is imperfective, and its aspect-partner \ru{прочита́ть} \emph{to read} is perfective, the two verbs having the same lexical (i.e. dictionary) meaning \emph{to read}; and again, the imperfective verb \ru{дели́ть} and its perfective partner \ru{раздели́ть} both mean \emph{to divide}. For many verbs, the perfective is inflected exactly like its imperfective partner (but with a prepositional prefix \ru{про-, раз-} etc.) except that, since a perfective verb cannot have a present tense, the forms corresponding to the imperfective present have a future meaning (the future perfective) e.g. \ru{чита́ем} \emph{we read}, \ru{прочита́ем} \emph{we shall read}. (For the future imperfective -- with no corresponding perfective forms -- see §14).

There are several possibilities for the ending of the first principal part of a verb (the infinitive), i.e. not only \ru{-ать} and \ru{-ить} as in \ru{чита́ть} and \ru{дели́ть}, but also e.g. \ru{-ять} as in \ru{меня́ть} \emph{to change} and, though much more rarely, \ru{-еть} as in \ru{име́ть} \emph{to have}, \ru{-ыть} as in \ru{откры́ть} \emph{to open}, \ru{-сть} as in \ru{класть} \emph{to lay}, \ru{-чь} as in \ru{мочь} \emph{to be able}, and one or two further endings.

But for the second principal part (the 1st person plural) there are only two possibilities:
\begin{enumerate}[label=\roman*)]
\item \ru{-ем} as in \ru{чита́ем} \emph{we read}, in which case the verb is said to belong to the first conjugation and its present tense (future for a perfective verb) is
\begin{center}\small
\begin{tabular}{@{}>{\cyrillicfont}ll >{\cyrillicfont}ll@{}}
чита́ет & it reads & прочита́ет & it will read\\
чита́ем & we read & прочита́ем & we shall read\\
чита́ют & they read & прочита́ют & they will read
\end{tabular}\end{center}
\item \ru{-им}, in which case the verb belongs to the second conjugation, and its present tense (future for a perfective verb) is
\begin{center}\small
\begin{tabular}{@{}>{\cyrillicfont}ll >{\cyrillicfont}ll@{}}
де́лит & it divides & раздели́т & it will divide\\
де́лим & we divide & раздели́м & we shall divide\\
де́лят & they divide & разделя́т & they will divide
\end{tabular}\end{center}
\end{enumerate}

All regular verbs with infinitives ending in \ru{-ить} belong to the second conjugation (with some monosyllabic exceptions of no mathematical interest) and all other regular verbs to the first, except (chiefly) the following verbs and their compounds:
\begin{center}\small
\begin{tabular}{@{}>{\cyrillicfont}ll@{}}
держа́ть, де́ржит & to hold, it holds\\
стоя́ть, стои́т & to stand, it stands\\
смотре́ть, смо́трит & to look at, it looks at\\
ви́деть, ви́дит & to see, it sees\\
висе́ть, виси́т & to hang, it hangs
\end{tabular}\end{center}

Verbs in \ru{-овать} (common in mathematics) change the \ru{-ов-} to \ru{-у-} in all forms derived from the second principal part; e.g. \ru{интегри́ровать} \emph{to integrate}, \ru{интегри́рует} \emph{it integrates}, \ru{интегри́руем} \emph{we integrate}, \ru{интегри́руют} \emph{they integrate}; similarly \ru{сле́довать} \emph{to follow}, \ru{сле́дует} \emph{it follows} etc.

Typical for the few verbs with infinitive ending \ru{-чь} is the present tense of \ru{мочь} \emph{to be able}, based on the root \ru{мог-} (with consonant variation): \ru{мо́жет} \emph{it can, it is possible}, \ru{мо́жем} \emph{we can}, \ru{мо́гут} \emph{they can}.

A transitive Russian verb can be made reflexive by suffixing \ru{-ся}, a short form of the reflexive \ru{себя́}, which is abbreviated to \ru{-сь} if the verb-form ends in a vowel (except that \ru{-ся} is used for all forms of the active participles; see §15). The resulting reflexive form will often be translated by a passive in English, e.g. \ru{чита́ется} \emph{it is being read} (lit. it is reading itself), \ru{де́лится} \emph{it is being divided}.

In modern Russian the present tense of the (irregular) verb \ru{быть} \emph{to be} has only the two forms \ru{есть} \emph{it is}, \ru{суть} \emph{they are}, and even these forms are usually omitted, their place being taken by a dash (or by nothing at all) or else by some circumlocution like \ru{явля́ется} \emph{it shows itself}, or \ru{представля́ет собо́й} \emph{it represents by means of itself}. The phrase \emph{there is not} is translated by the single word \ru{нет}, an abbreviation of \ru{не есть}.

\section{The future imperfective and the past tense}

The future of the irregular imperfective verb \ru{быть} \emph{to be} is
\begin{center}\small
\begin{tabular}{@{}>{\cyrillicfont}ll@{}}
бу́дет & it will be\\
бу́дем & we shall be\\
бу́дут & they will be
\end{tabular}\end{center}
and the future of any other imperfective verb is formed by combining its infinitive with the future of this auxiliary verb \ru{быть}:
\begin{center}\small
\begin{tabular}{@{}>{\cyrillicfont}ll >{\cyrillicfont}ll@{}}
бу́дет чита́ть & it will read & бу́дет дели́ть & it will divide\\
бу́дем чита́ть & we shall read & бу́дем дели́ть & we shall divide\\
бу́дут чита́ть & they will read & бу́дут дели́ть & they will divide
\end{tabular}\end{center}

For the regular verbs illustrated here, the past tense (perfective or imperfective) is formed by replacing the infinitive \ru{-ть} with \ru{-л, -ла, -ло, -ли}; e.g. \ru{чита́л, чита́ла, чита́ло; чита́ли} (perf. \ru{прочита́л}, etc.) \emph{it read, or has read; we or they read, or have read}, and similarly \ru{дели́л, дели́ла, дели́ло; дели́ли} (or \ru{раздели́л}, etc.) \emph{it divided}, and from \ru{быть} \emph{to be}: \ru{был, была́, бы́ло; бы́ли}.\corrnote{原文把“去掉不定式 \ru{-ть} 加 \ru{-л/-ла/-ло/-ли}”写成所有动词的通则；对规则例词成立，但 \ru{мочь}\(\to\)\ru{мог}, \ru{нести}\(\to\)\ru{нёс}, \ru{идти}\(\to\)\ru{шёл} 等需另行处理，故正文把范围限定为“the regular verbs illustrated here”。}

These past active forms, e.g. \ru{чита́л} referring to a masculine singular, \ru{чита́ла} to a feminine singular, \ru{чита́ло} to a neuter singular, \ru{чита́ли} to any plural, vary for gender like the short predicate adjectives (see §8, Remark iv); the explanation being that in the Slavic languages all verbs form their compound past tenses with \emph{to be} (instead of \emph{to have} as in English, e.g. \emph{they have read}; but compare \emph{they are gone} in English), and in Russian (as usual) the forms of the verb \emph{to be} have been omitted, so that the final result is, for example, \ru{они́ чита́ли} \emph{they read} (past tense).

Note that \ru{что́бы} \emph{in order that} (p. 50) is followed by a past tense to express purpose: \ru{что́бы отноше́ние бы́ло равно́} \emph{in order that the ratio may (might) be equal}.

Here the particle \ru{бы}, an old conditional form from the verb \ru{быть} \emph{to be}, is amalgamated with \ru{что} \emph{that} to express purpose. In conditional sentences expressing a supposition contrary-to-fact, it is used as a separate word, again with the past tense of the verb, in both the if-clause and the then-clause: \ru{е́сли бы число́ p бы́ло соста́вным, то оно́ бы́ло бы разложи́мым} \emph{if the number p were composite, then it would be factorable}.

\begin{keybox}[Exercise 3.8]
Read aloud the following sentences, illustrating non-participial verb forms.
\end{keybox}

\begin{enumerate}[label=\arabic*)]
\item \ru{рассма́триваем систе́му координа́т}\\
      we consider a system of coordinates.
\item \ru{положе́ние то́чки M определя́ется угло́м $\theta$}\\
      the position of the point M is determined by the angle $\theta$.
\item \ru{уравне́ние представля́ет собо́й пряму́ю}\\
      the equation represents a straight line.
\item \ru{скаля́рное произведе́ние обраща́ется в нуль}\\
      the scalar product reduces to zero.
\item \ru{э́то усло́вие приво́дит к уравне́нию}\\
      this condition leads to the equation.
\item \ru{мы прихо́дим к теоре́ме}\\
      we arrive at the theorem.
\item \ru{об эквивале́нтных мно́жествах говоря́т, что они́ име́ют одина́ковую мо́щность}\\
      about equivalent sets they say that they have the same cardinality.\\[-1mm]
      \emph{(Compare the impersonal/indefinite-personal ``they say'', similarly ``they write'', ``they define''.)}
\item \ru{по́ле не име́ет дели́телей нуля́}\\
      a field does not have divisors of zero.
\item \ru{ка́ждое рациона́льное число́ принадлежи́т одному́ из двух мно́жеств}\\
      every rational number belongs to one of the two sets.
\item \ru{име́ется тако́е мно́жество D, кото́рое \ldots}\\
      there is such a set D, which \ldots
\item \ru{е́сли мы занумеру́ем рациона́льные то́чки \ldots}\\
      if we enumerate the rational points \ldots
\end{enumerate}

The following sentences illustrate the important predicative use of the instrumental case:

\begin{enumerate}[label=\arabic*),start=12]
\item \ru{чи́сла $\rho$ и $\theta$ называ́ются поля́рными координа́тами}\\
      the numbers $\rho$ and $\theta$ are called polar coordinates.
\item \ru{э́тот необходи́мый при́знак не явля́ется доста́точным}\\
      this necessary criterion is not sufficient.
\item \ru{число́ 1 не счита́ется ни просты́м, ни соста́вным}\\
      the number 1 is considered neither prime nor composite.
\end{enumerate}

The predicate instrumental is particularly common with verbs such as \ru{называ́ть(ся)} ``call/be called'', \ru{явля́ться} ``be, constitute'', and \ru{счита́ть(ся)} ``consider/be considered''.

\section{The four adjectival and two adverbial participles}

The first four of the following six rules for the formation of participles refer to the four adjectival participles: present active, past active, present passive and past passive, which are verbal adjectives declined like \ru{ра́вный} or \ru{после́дний} in §8. But rules 5, 6 refer to the two adverbial participles (present and past), called adverbial because they are indeclinable substitutes (though formerly they too were declined) for the two active adjectival participles. When the word ``participle'' is used alone, it usually refers to an adjectival participle.

\paragraph{Rule 1.} In the present active participle (formed only for imperfective verbs) the \ru{-т} of the third person plural indicative (e.g. \ru{чита́ют} \emph{they read}) is replaced by \ru{-щий}: \ru{чита́ют} \(\to\) \ru{чита́ющий} \emph{reading}; \ru{де́лят} \(\to\) \ru{деля́щий} \emph{dividing}.

\paragraph{Rule 2.} In the past active participle the \ru{-ть} of the infinitive is replaced, for the regular examples here, by \ru{-вший}: \ru{чита́ть} \(\to\) \ru{чита́вший} \emph{having read}, \ru{дели́ть} \(\to\) \ru{дели́вший} \emph{having divided}.\corrnote{原文把这一替换写成一般规则；辅音词干等还可用 \ru{-ш-}，故只把适用范围限定为本书所示规则例词。}

\paragraph{Rule 3.} In the present passive participle (only for imperfective verbs) the ending \ru{-мый} is added to the first person plural: \ru{чита́ем} \emph{we read} \(\to\) \ru{чита́емый} \emph{being read}; \ru{де́лим} \emph{we divide} \(\to\) \ru{дели́мый} \emph{being divided}. This participle often carries the meaning of \emph{-able} or \emph{-ible}; thus \ru{дели́мый} \emph{divisible}.

\paragraph{Rule 4.} In the past passive participle the \ru{-ть} of the infinitive is replaced by \ru{-енный} (for infinitives in \ru{-ить}, the \ru{-ить} is replaced by \ru{-енный} or \ru{-ённый}): \ru{чита́ть} \(\to\) \ru{чи́танный} \emph{having been read}; \ru{дели́ть} \(\to\) \ru{разделённый} \emph{having been divided}; but some verbs have \ru{-тый}, e.g. \ru{закры́ть} \(\to\) \ru{закры́тый} \emph{closed}.

The short predicate-adjective form (see §8) is especially common for the participles under Rules 3 and 4; e.g. \ru{разложи́м} \emph{factorable} (\ru{неразложи́м} \emph{nonfactorable}) from \ru{разложи́ть} \emph{to decompose, factor}; \ru{дан} for \ru{да́нный} as in \ru{пусть дан алгори́тм} \emph{let there be given an algorithm}, \ru{закры́т} \emph{closed}, etc.

\paragraph{Rule 5.} In the present adverbial participle (usually only for imperfective verbs) the ending of the third plural is replaced by \ru{-я}: \ru{чита́ют} \(\to\) \ru{чита́я} \emph{reading}; \ru{де́лят} \(\to\) \ru{деля́} \emph{dividing}.

\paragraph{Rule 6.} In the past adverbial participle (usually only for perfective verbs) the \ru{-ть} of the infinitive is replaced by \ru{-в} (less often \ru{-вши}): \ru{прочита́ть} \(\to\) \ru{прочита́в} or \ru{прочита́вши} \emph{having read}; \ru{раздели́ть} \(\to\) \ru{раздели́в} or \ru{раздели́вши} \emph{having divided}.

\begin{keybox}[Exercise 3.9]
Read the following six sets of sentences illustrating the six participial forms, four adjectival, and two adverbial. Since participles are extremely important in mathematical Russian, the endings illustrated here, and in particular the agreement of participles with nouns, will repay careful attention. Note that a postposed expanded participial phrase is normally set off by a comma, which must often be omitted in English.\corrnote{原文绝对写作 “if the participle follows the noun, there is a comma”。俄语标点规则更具体：典型的是后置的扩展分词短语用逗号隔开；单个、限制性或词汇化结构并不服从这一绝对规则，故作最小收窄。}
\end{keybox}


\begingroup\small\sloppy\subsection*{1. Present active participle: \ru{-ющий/-ящий}}

\begin{enumerate}[label=1.\arabic*)]
\item \ru{пусть M --- произво́льная то́чка прямо́й, проходя́щей че́рез A}\\
      let M be an arbitrary point of the line passing through A.\\[-1mm]
      \emph{(Here \ru{проходя́щей} is genitive feminine singular, agreeing with \ru{прямо́й}.)}
\item \ru{пересече́нием мно́жеств A и B называ́ется мно́жество элеме́нтов, принадлежа́щих и A и B}\\
      by the intersection of the sets A and B is meant the set of elements belonging both to A and to B.
\item \ru{взаи́мно однозна́чным соотве́тствием называ́ется соотве́тствие, име́ющее сле́дующие сво́йства}\\
      by a one-to-one correspondence is meant a correspondence having the following properties.
\item \ru{существу́ет элеме́нт, лежа́щий ме́жду $a$ и $b$}\\
      there exists an element lying between $a$ and $b$.
\item \ru{опера́ции, интересу́ющие нас \ldots}\\
      the operations interesting us \ldots
\item \ru{элеме́нт $b$ из S, не явля́ющийся алгебраи́ческим над R, называ́ется трансценде́нтным над R}\\
      an element $b$ from S not algebraic over R is called transcendental over R.
\item \ru{реше́ния уравне́ния, удовлетворя́ющие усло́виям \ldots}\\
      the solutions of the equation satisfying the conditions \ldots
\item \ru{относи́тельно систе́м фу́нкций, образу́ющих класс K}\\
      with respect to the systems of functions forming the class K.
\item \ru{рассмо́трим сумми́рование расходя́щихся рядо́в}\\
      we shall consider the summation of divergent series.
\end{enumerate}

\subsection*{2. Past active participle: \ru{-вший}}

\begin{enumerate}[label=2.\arabic*)]
\item \ru{перечи́слим не́которые кла́ссы фу́нкций, получи́вших назва́ние элемента́рных}\\
      we shall enumerate some classes of functions having received the name of elementary.
\item \ru{э́та теоре́ма принадлежи́т Лебе́гу, сформули́ровавшему её в 1911 г.}\\
      this theorem is due to Lebesgue, who formulated it in 1911.
\end{enumerate}

\subsection*{3. Present passive participle: \ru{-мый}}

\begin{enumerate}[label=3.\arabic*)]
\item \ru{полу́чим иско́мое уравне́ние}\\
      we obtain the desired equation.
\item \ru{величина́ отре́зка, отсека́емого на о́си $Ox$}\\
      the magnitude of the segment cut off on the axis $Ox$.
\item \ru{фу́нкция, представля́емая многочле́ном}\\
      a function represented (or representable) by a polynomial.
\item \ru{фу́нкция называ́ется дифференци́руемой}\\
      a function is called differentiable.
\item \ru{так называ́емые логи́ческие переме́нные}\\
      the so-called logical variables.
\item \ru{соотве́тствие, устана́вливаемое таки́м о́бразом}\\
      the correspondence established in such a manner.
\item \ru{э́ти фу́нкции облада́ют интегри́руемыми квадра́тами}\\
      these functions possess integrable squares.
\item \ru{вся́кий ряд, сумми́руемый ме́тодом Чеза́ро}\\
      every series summable by the method of Cesàro.
\item \ru{обобще́ние рима́новой геоме́трии, применя́емое в о́бщей тео́рии относи́тельности}\\
      the generalization of Riemannian geometry applied in the general theory of relativity.
\item \ru{предполага́ются вы́полненными сле́дующие усло́вия, называ́емые аксио́мами}\\
      the following conditions, called axioms, are assumed fulfilled.
\item \ru{топологи́ческое простра́нство, гомеомо́рфное метри́ческому простра́нству, называ́ется метризу́емым простра́нством}\\
      a topological space homeomorphic to a metric space is called a metrizable space.
\end{enumerate}

\subsection*{4. Past passive participle: \ru{-енный/-анный/-тый}}

\begin{enumerate}[label=4.\arabic*)]
\item \ru{отноше́ние, равно́ за́данному числу́}\\
      a ratio equal to a given number.
\item \ru{пусть даны́ то́чки $M_1$ и $M_2$}\\
      let there be given points $M_1$ and $M_2$.
\item \ru{теоре́ма дока́зана}\\
      the theorem is proved.
\item \ru{значе́нию $x$ при́дано прираще́ние $\Delta x$}\\
      to the value of $x$ there is given an increment $\Delta x$.
\item \ru{из дока́занной фо́рмулы}\\
      from the proved formula.
\item \ru{обобщённая теоре́ма о сре́днем}\\
      the generalized theorem about the mean.
\item \ru{э́то ра́венство, перепи́санное в ви́де \ldots}\\
      this equality, rewritten in the form \ldots
\item \ru{рассмо́трим примене́ние, свя́занное с интегра́льным исчисле́нием}\\
      we shall consider an application connected with the integral calculus.
\item \ru{э́та фу́нкция не мо́жет быть изображена́ одни́м перехо́дом к преде́лу}\\
      this function cannot be represented by one passage to the limit.
\item \ru{рассмо́тренная зада́ча поставлена́ некорре́ктно}\\
      the considered problem is posed incorrectly.
\item \ru{мы изло́жим не́которые вопро́сы тео́рии обобщённых фу́нкций, постро́енной С. Л. Со́болевым и Л. Шва́рцем}\\
      we shall explain some questions of the theory of generalized functions constructed by S. L. Sobolev and L. Schwartz.
\item \ru{мно́жество F, образо́ванное опи́санным о́бразом \ldots}\\
      the set F formed in the described manner \ldots
\end{enumerate}

\subsection*{5. Imperfective adverbial participle: \ru{-ая/-я}}

\begin{enumerate}[label=5.\arabic*)]
\item \ru{обозна́чим че́рез $(\rho,\theta)$ поля́рные координа́ты то́чки M, счита́я поля́рной о́сью $Ox$}\\
      we shall denote by $(\rho,\theta)$ the polar coordinates of the point M, considering $Ox$ as polar axis.
\item \ru{вычита́я из ра́венства (1) ра́венство (2), полу́чим иско́мое уравне́ние}\\
      subtracting equality (2) from equality (1), we obtain the sought equation.
\item \ru{переходя́ от поля́рных координа́т к декартовым}\\
      passing from polar coordinates to Cartesian coordinates.
\item \ru{обознача́я э́то число́ че́рез $-D$, полу́чим \ldots}\\
      denoting this number by $-D$, we obtain \ldots
\item \ru{сле́дуя Дедеки́нду}\\
      following Dedekind.
\item \ru{придава́я аргуме́нту $x$ прираще́ние $\Delta x$}\\
      giving the argument $x$ the increment $\Delta x$.
\item \ru{полага́я $n=1$ в теоре́ме Те́йлора}\\
      setting $n=1$ in Taylor's theorem.
\item \ru{подставля́я в э́ту аксио́му фо́рмулу A, полу́чим \ldots}\\
      substituting formula A into this axiom, we obtain \ldots
\item \ru{дифференци́руя под зна́ком интегра́ла}\\
      differentiating under the integral sign.
\item \ru{производя́ интегри́рование, мо́жно написа́ть}\\
      carrying through the integration, it is possible to write.
\item \ru{зна́я пе́рвую квадрати́чную фо́рму, мо́жно измеря́ть длину́}\\
      knowing the first quadratic form, it is possible to measure the length.
\end{enumerate}

(Note that in standard Russian an adverbial participle normally has the same understood subject as the finite predicate; constructions violating this requirement can be felt as dangling.)\corrnote{原文此处写作 “in Russian the adverbial participles are never felt to be dangling.” 这与现代规范不符；\ru{деепричастный оборот} 的动作执行者通常须与主句谓语的主语一致，故作最小更正。}

\subsection*{6. Perfective adverbial participle: \ru{-в/-вши}}

\begin{enumerate}[label=6.\arabic*)]
\item \ru{раздели́в на $g$, получа́ем \ldots}\\
      dividing by $g$ we obtain \ldots
\item \ru{постоя́нная C определя́ется, положи́в $x=a$}\\
      the constant C is determined by setting $x=a$.
\item \ru{результа́т не мо́жно уточни́ть, замени́в $1/b^2+\varepsilon$ че́рез $1/b^2$}\\
      the result cannot be sharpened by replacing $1/b^2+\varepsilon$ by $1/b^2$.
\item \ru{подста́вив э́ти выраже́ния в интегра́л}\\
      substituting these expressions into the integral.
\end{enumerate}
\endgroup
\fussy